\chapter{Épreuve chronométrée 9}\label{ch:iv1:timed-test-9}

\begin{ivtrains}
Il faut la traiter d'une seule traite, sans notes ni calculatrice.
\end{ivtrains}

\ivtestsection{A}{6}{40}

\begin{question}\label{iq:iv1:timed-test-9:1}
Deux pièces équilibrées sont lancées ensemble, tour après tour. Une pièce ayant déjà montré face peut être mise de côté. Quelle est l'espérance du nombre de tours nécessaires pour que chaque pièce ait montré face au moins une fois~?
\end{question}

\begin{question}\label{iq:iv1:timed-test-9:2}
Trois sommets distincts d'un cube sont choisis au hasard. Quelle est la probabilité qu'ils forment un triangle équilatéral~?
\end{question}

\begin{question}\label{iq:iv1:timed-test-9:3}
Un paramètre a pour loi a priori $N(0, 1)$. Vous observez une seule mesure égale à 2, avec un bruit gaussien indépendant de variance 1. Quelles sont la moyenne et la variance a posteriori~? Et si la variance du bruit vaut 4~?
\end{question}

\begin{question}\label{iq:iv1:timed-test-9:4}
Un nombre est tiré uniformément parmi les diviseurs de $10^6$. Quelle est la probabilité qu'il soit un multiple de $10^3$~?
\end{question}

\begin{question}\label{iq:iv1:timed-test-9:5}
$X$ et $Y$ sont des normales centrées réduites indépendantes. Quelle est la probabilité que $X > Y + 1$, sachant que $P(Z < 0.71) = 0.76$ pour une normale centrée réduite $Z$~? Quelle est $\mathbb{E}[X \mid X + Y = 2]$~?
\end{question}

\begin{question}\label{iq:iv1:timed-test-9:6}
$X$ et $Y$ sont deux variables exponentielles indépendantes de moyenne 1. Quelles sont les moyennes du plus petit et du plus grand des deux~? Quelle est la probabilité que $X > Y + 1$~?
\end{question}

\ivtestsection{B}{6}{60}

\begin{question}\label{iq:iv1:timed-test-9:7}
$N$ suit une loi de Poisson de moyenne 1. Quelle est la probabilité que $N$ soit pair~? Par ailleurs, des clients arrivent dans une boutique selon un processus de Poisson au taux de 6 par heure, et chacun achète avec probabilité $\frac12$, indépendamment des autres. Quelle est la probabilité qu'il n'y ait aucune vente en une heure~?
\end{question}

\begin{question}\label{iq:iv1:timed-test-9:8}
Cinq dés équilibrés sont lancés. Quelle est la probabilité d'obtenir exactement deux paires, c'est-à-dire deux valeurs apparaissant chacune deux fois et une troisième valeur apparaissant une seule fois~?
\end{question}

\begin{question}\label{iq:iv1:timed-test-9:9}
Deux actifs non corrélés ont des rendements excédentaires espérés de 5~\% et 10~\%, et des volatilités de 10~\% et 20~\%. Quel portefeuille maximise le ratio de Sharpe, et quelle est la valeur de ce ratio~?
\end{question}

\begin{question}\label{iq:iv1:timed-test-9:10}
Une main de cinq cartes est distribuée à partir d'un jeu classique de 52 cartes, qui compte $2\,598\,960$ mains possibles. Quelle est la probabilité d'obtenir exactement deux paires, c'est-à-dire deux cartes d'une valeur, deux cartes d'une autre valeur et une carte d'une troisième valeur~?
\end{question}

\begin{question}\label{iq:iv1:timed-test-9:11}
Un avion compte 10 sièges et 10 passagers, chacun ayant un siège assigné. Le premier a perdu sa carte d'embarquement et s'assoit au hasard. Chaque passager suivant prend son propre siège s'il est libre, et sinon un siège libre au hasard. Quelle est la probabilité que le dernier passager obtienne son siège~? Quelle est la probabilité que le neuvième passager ne l'obtienne pas~?
\end{question}

\begin{question}\label{iq:iv1:timed-test-9:12}
En utilisant l'approximation de Stirling $\ln n! \approx n\ln n - n + \frac12\ln(2\pi n)$, combien de chiffres décimaux possède $100!$~? Vous pouvez utiliser $\ln 10 = 2.303$. Vers quelle limite tend $(n!)^{1/n}/n$ lorsque $n$ croît~?
\end{question}
